\lesson{5}{Oct 18 2021 Mon (11:31:19)}{Completing the Square}{Unit 2}

\subsubsection*{Intro to the Discriminant}

The formula for the \bf{discriminant} is:

\begin{align}
    b^2 - 4ac
\end{align}

Let's try a quick example:

\begin{example}
    \begin{align}
        f(x) &= x^2 + 12x + 26 \\
             &= b^2 - 4ac \\
             &= (12)^2 - 4(1)(26) \\
             &= 144 - 104 \\
             &= 40
    \end{align}
    
    So, the discriminant isn't a perfect square. This tells us that there will be two irrational numbers.
\end{example}

\subsubsection*{Using Vertex Form: Minimum or Maximum}

\begin{example}
    \begin{align}
        f(x) = (x + 6)^2 - 10
    \end{align}
    
    From this equation, we get the vertex of the parabola, which is $(-6, -10)$
    But, it raises the question. Is it the vertex minimum, or the vertex maximum?
    
    To solve this, you must look at the value or the leading coefficient to make this determination.
    
    In this case, it would be $a = 1$, which is positive, meaning it goes upward.
\end{example}

\newpage
